\documentclass[11pt]{article}
\newcommand{\ShortTitle}{CATALAN'S CONSTANT}
\usepackage[T1]{fontenc}
\usepackage[utf8]{inputenc}
\usepackage{mathpazo}
\usepackage[scaled=0.92]{helvet}
\usepackage{courier}
\usepackage{amsmath,amssymb,mathtools}
\usepackage[a4paper,landscape,left=18mm,right=18mm,top=19mm,bottom=17mm,headheight=14pt]{geometry}
\usepackage{graphicx,microtype,xcolor,fancyhdr,hyperref,booktabs,array}
\definecolor{Ink}{HTML}{243138}
\definecolor{Teal}{HTML}{21676B}
\definecolor{Copper}{HTML}{A45C36}
\definecolor{Soft}{HTML}{E9F0EF}
\definecolor{Gray}{HTML}{627078}
\hypersetup{colorlinks=true,linkcolor=Teal,citecolor=Teal,urlcolor=Teal,pdfauthor={Komorebi P.}}
\urlstyle{same}
\setlength{\parindent}{0pt}
\setlength{\parskip}{0.55em}
\setlength{\emergencystretch}{1.5em}
\linespread{1.04}
\pagestyle{fancy}
\fancyhf{}
\fancyhead[L]{\sffamily\scriptsize\color{Gray}OPERATION KOMOREBI / VISUAL GUIDE}
\fancyhead[R]{\sffamily\scriptsize\color{Gray}\ShortTitle}
\fancyfoot[L]{\sffamily\scriptsize\color{Gray}READER EDITION 1.0 / 9 OCTOBER 2026}
\fancyfoot[R]{\sffamily\small\color{Ink}\thepage}
\renewcommand{\headrulewidth}{0.25pt}
\renewcommand{\footrulewidth}{0pt}
\providecommand{\N}{\mathbb N}
\providecommand{\Z}{\mathbb Z}
\providecommand{\Q}{\mathbb Q}
\providecommand{\R}{\mathbb R}
\providecommand{\C}{\mathbb C}
\providecommand{\eps}{\varepsilon}
\providecommand{\dd}{\mathrm d}
\providecommand{\norm}[1]{\left\lVert #1\right\rVert}
\providecommand{\abs}[1]{\left\lvert #1\right\rvert}
\newcommand{\GuideCover}[4]{%
\hypersetup{pageanchor=false}
\begin{titlepage}\thispagestyle{empty}
{\sffamily\small\color{Teal}OPERATION KOMOREBI\hfill CONJECTURE SERIES #1 / VISUAL GUIDE}\par
\vspace{18mm}
{\fontsize{40}{44}\selectfont\color{Ink}#2}\par
\vspace{4mm}{\LARGE\color{Ink}#3}\par
\vspace{6mm}{\color{Teal}\rule{\linewidth}{1.1pt}}\par
\vspace{4mm}{\large Komorebi P.}\par
{\sffamily\small\color{Gray}READER EDITION 1.0 / 9 OCTOBER 2026}\par
\vfill
\begin{minipage}{0.78\linewidth}\large #4\end{minipage}\par
\vspace{5mm}{\sffamily\small\color{Teal}WORKED EXAMPLES\quad /\quad EXACT COORDINATES\quad /\quad OPEN PROBLEMS}\par
\end{titlepage}\hypersetup{pageanchor=true}}
\newcommand{\GuideSpread}[5]{%
\clearpage
{\fontsize{22}{26}\selectfont\bfseries\color{Ink}#1}\par
\vspace{1mm}{\sffamily\small\color{Gray}#2}\par
\vspace{3mm}
\noindent\begin{minipage}[t]{0.665\linewidth}\vspace{0pt}
\centering\includegraphics[width=\linewidth,height=113mm,keepaspectratio]{#3}
\end{minipage}\hfill
\begin{minipage}[t]{0.295\linewidth}\vspace{0pt}
{\sffamily\small\bfseries\color{Teal}READ THE FIGURE}\par
\small #4
\end{minipage}\par
\vfill
\noindent\colorbox{Soft}{\parbox{\dimexpr\linewidth-2\fboxsep\relax}{\vspace{1.3mm}\small\textbf{Carry forward.} #5\vspace{1.3mm}}}\par
}

\hypersetup{pdftitle={Catalan's constant: a visual guide},pdfsubject={Conjecture Series 08. Exact identities, finite evidence and the unresolved irrationality question.}}
\begin{document}
\GuideCover{08}{Catalan's constant}{A visual guide to the irrationality question}{The constant is easy to define and difficult to classify. Follow the exact identities, the rational approximants and the arithmetic cost of turning a small real error into a small integer linear form. The guide marks where the finite evidence ends and the open problem begins.}

\GuideSpread{A number we can locate}{01 / THE QUESTION}{figures/01_constant.pdf}{
Catalan's constant is
\[
G=\sum_{k=0}^{\infty}\frac{(-1)^k}{(2k+1)^2}.
\]
The series converges. It defines a particular real number, about $0.9159655942$.

The conjectural assertion is $G\notin\Q$. Its decimal expansion, however many digits we compute, cannot decide that assertion.

This edition records research and audited computations. It does not establish irrationality.
}{A definition and an accurate decimal value leave the rationality question open.}

\GuideSpread{Alternating terms bracket the constant}{02 / AN ELEMENTARY CERTIFIED INTERVAL}{figures/02_alternating_series.pdf}{
Let $s_N$ contain the first $N$ terms. The odd partial sums lie above $G$. The even partial sums lie below it.

The next omitted term gives
\[
|G-s_N|\leq\frac{1}{(2N+1)^2}.
\]
Every plotted partial sum is an exact rational number. The horizontal marker for $G$ is a decimal display.

The bound certifies approximation. It gives no information about whether $G$ is rational.
}{The series supplies exact rational enclosures with a transparent error bound.}

\GuideSpread{The same constant as an area}{03 / A SINGLE INTEGRAL}{figures/03_log_integral.pdf}{
Expanding $(1+x^2)^{-1}$ and integrating the terms gives
\[
G=\int_0^1\frac{-\log x}{1+x^2}\,\dd x.
\]
The logarithm diverges at $0$. Its integral over a short interval remains finite.

The graph clips the visible height at $5.5$. The shaded display is not a numerical proof of the identity.

Integral identities can expose positivity and estimates that a decimal expansion hides.
}{An exact integral representation changes the available tools without resolving rationality.}

\GuideSpread{A square instead of a line}{04 / A DOUBLE INTEGRAL}{figures/04_double_integral.pdf}{
On the unit square,
\[
G=\int_0^1\!\int_0^1
\frac{\dd x\,\dd y}{1+x^2y^2}.
\]
The integrand stays between $1/2$ and $1$.

Integrating each geometric term produces $(2k+1)^{-2}$. Absolute convergence of the integrated terms justifies that passage.

The contour colors show the integrand. They are floating display coordinates, not evidence from quadrature.
}{The squared odd denominators arise from integrating once in each coordinate.}

\GuideSpread{A family with a difficult middle value}{05 / DIRICHLET BETA}{figures/05_beta_family.pdf}{
The beta function packages these series as
\[
\beta(s)=\sum_{k\geq0}\frac{(-1)^k}{(2k+1)^s}
\quad(s>0).
\]
At positive odd integers, familiar evaluations involve powers of $\pi$.

Catalan's constant is the positive even value $\beta(2)$. The neighbouring formulas do not yield a comparable evaluation for it.

The checked sources establish neither its irrationality nor its transcendence. A September 2026 proof claim and a subsequent critical note remain a source dispute, recorded in the monograph.
}{A successful formula at $s=1$ or $s=3$ does not settle the value at $s=2$.}

\GuideSpread{Why small integer forms would settle the question}{06 / THE LATTICE GAP}{figures/06_rationality_lattice.pdf}{
Suppose $G=a/b$ with integers $a,b$ and $b>0$.

For integer $A,B$,
\[
b(AG-B)=aA-bB\in\Z.
\]
If the form is nonzero, its absolute value is at least $1/b$.

Thus an infinite sequence of nonzero integer forms tending to zero contradicts the assumed rationality.

All three words matter: integer, nonzero and vanishing.
}{The route is to force a nonzero integer into a gap too small to contain one.}

\GuideSpread{A good decimal approximation has a denominator}{07 / THE FIRST ARITHMETIC OBSTACLE}{figures/07_series_denominators.pdf}{
Write $s_N=P_N/D_N$ in lowest terms.

The left panel plots exact reduced denominators through $N=60$. The right panel plots the guaranteed alternating error bound.

These quantities move on very different scales. Multiplying the approximation error by $D_N$ can erase the advantage of the approximation.

The displayed error bound is an upper bound. Its failure to prove a small linear form does not prove that no sharper bound exists.
}{The relevant quantity is the error after the rational denominators are cleared.}

\GuideSpread{Clearing is a multiplication, with a cost}{08 / FROM A RATIONAL APPROXIMANT TO A LINEAR FORM}{figures/08_error_clearing.pdf}{
The identity is exact:
\[
D_NG-P_N=D_N(G-s_N).
\]
It produces integer coefficients. It also multiplies the error by the reduced denominator.

Fast approximation and controlled denominators must work together. A construction needs an estimate for their product.

Later common divisors can reduce that product. A chosen scale must be accounted for when measuring those divisors.
}{Track analytic decay and arithmetic growth in the same normalization.}

\GuideSpread{A recurrence creates two coordinated sequences}{09 / THE ZUDILIN RECURRENCE}{figures/09_recurrence_structure.pdf}{
Set $p(n)=20n^2-8n+1$ and
\begin{align*}
A_n&=(2n+1)^2(2n+2)^2p(n),\\
C_n&=(2n-1)^2(2n)^2p(n+1).
\end{align*}
The middle coefficient is
\begin{align*}
Q_n={}&3520n^6+5632n^5\\
&+2064n^4-384n^3\\
&-156n^2+16n+7.
\end{align*}
Both sequences satisfy the displayed relation for $n\geq1$. Their different initial values make their roles distinct.
}{The recurrence is exact algebra. Approximation and denominator claims require additional arguments.}

\GuideSpread{Start with exact rational values}{10 / INITIAL DATA FIX THE NORMALIZATION}{figures/10_initial_sequences.pdf}{
The historical audit uses
\[
u_0=1,\quad u_1=7/4,
\qquad v_0=0,\quad v_1=13/8.
\]
With these initial values, the recurrence determines every later rational term.

The quotient $v_n/u_n$ is the approximant in this normalization. A source using an integral pair may multiply both sequences by a signed constant.

The source bundle records that translation explicitly. It must be applied before comparing equations or denominators.
}{Initial values are part of the mathematical definition, not an editorial detail.}

\GuideSpread{The characteristic equation explains the scale}{11 / GROWTH IS A SOURCE DEPENDENT CLAIM}{figures/11_recurrence_growth.pdf}{
The leading polynomial coefficients give
\[
\lambda^2-11\lambda-1=0,
\qquad
\lambda_{\pm}=\frac{11\pm5\sqrt5}{2}.
\]
The finite curve plots $\log_{10}(u_n)/n$ through $n=60$.

The dominant positive root predicts the growth scale. Turning that prediction into a theorem about a selected solution requires the appropriate recurrence analysis and initial conditions.

This plot is a consistency display. It is not that theorem.
}{Leading coefficients identify a candidate scale. A rigorous asymptotic statement needs its hypotheses.}

\GuideSpread{The approximants quickly approach the target}{12 / CERTIFIED FINITE ERRORS}{figures/12_approximant_error.pdf}{
Each point uses exact $u_n,v_n$ and a certified rational interval for $G$.

The interval determines the error's sign and brackets its magnitude. The plotted logarithm is a rounded display of that bracket.

The selected indices are $0,1,2,5,10,20,40$. Teal and copper distinguish the two error signs.

This is much faster approximation than the elementary partial sums. The rational denominators still have to be tracked.
}{A tiny certified error is useful evidence about an approximant, not an irrationality proof by itself.}

\GuideSpread{The determinant stays away from zero}{13 / A WRONSKIAN IDENTITY}{figures/13_wronskian.pdf}{
Define
\[
W_n=u_{n+1}v_n-u_nv_{n+1}.
\]
Its exact formula is
\[
W_n=\frac{(-1)^{n+1}p(n+1)}{2(2n+1)^2(2n+2)^2}.
\]
The numerator cannot vanish for $n\geq0$. Adjacent approximants are distinct.

The left panel shows the alternating sign. The right panel shows the decreasing magnitude on a logarithmic scale.

This determinant does not establish nonvanishing of an arbitrary filtered combination.
}{An exact determinant proves independence of the adjacent pairs. Its scope should remain explicit.}

\GuideSpread{Cancellation makes the determinant tractable}{14 / THE IDENTITY HAS AN ELEMENTARY PROOF}{figures/14_wronskian_identity.pdf}{
Substitute the recurrence into both terms of $W_n$.

The contributions containing $Q_n$ cancel. The remaining terms give
\[
W_n=-\frac{C_n}{A_n}W_{n-1}.
\]
The coefficient product telescopes. The initial determinant is $W_0=-13/8$.

This establishes the formula for every nonnegative index. The finite check through $n=200$ separately checks the implementation.
}{A universal written proof and a finite computational check are different kinds of evidence.}

\GuideSpread{Scale the first coefficient before filtering}{15 / INTEGER PACKET INPUTS}{figures/15_scaled_coefficients.pdf}{
The packet construction writes
\[
T_n=16^n u_n.
\]
The displayed values are exact integers. There is also a proof for every $n$.

Krattenthaler and Rivoal's coefficient formula gives
\[
u_n=\sum_{j=0}^n\binom nj
\binom{n-1/2}{j}\binom{n+j-1/2}{j}.
\]
For integer $h\geq0$, put $b_{h,j}=4^j\binom{h-1/2}{j}$. These are integer coefficients of
\[
(1+4z)^h(1+4z)^{-1/2},
\]
since the second factor has coefficients $(-1)^j\binom{2j}{j}$.

Each summand of $T_n$ is therefore $\binom nj16^{n-j}b_{n,j}b_{n+j,j}\in\Z$.

The separate finite audit checked every recurrence division through $n=10000$.
}{A known coefficient formula supplies a universal proof. The exact division scan separately checks the implementation.}

\GuideSpread{A binomial packet has a local update rule}{16 / PASCAL ADDITION}{figures/16_pascal_transform.pdf}{
Define
\[
S_{n,m}=\sum_{i=0}^m\binom mi T_{n+i}.
\]
Pascal's identity yields
\[
S_{n,m+1}=S_{n,m}+S_{n+1,m}.
\]
The initial row is $S_{n,0}=T_n$.

This local recurrence allows a large triangular scan without rebuilding every binomial sum. It is also valid modulo $2^K$.
}{The fast computation follows an exact binomial identity.}

\GuideSpread{What is being added in one packet}{17 / A WORKED FINITE EXAMPLE}{figures/17_binomial_packet.pdf}{
For $n=1,m=4$, the weights are
\[
1,\ 4,\ 6,\ 4,\ 1.
\]
They multiply five consecutive terms $T_1,\ldots,T_5$.

The bars show the positive summands on a logarithmic axis. The final sum is computed as an exact integer, with no floating summation.

Large summands do not directly determine the divisibility of their sum. Cancellation in residue classes is an arithmetic question.
}{Size and divisibility are separate properties of the packet sum.}

\GuideSpread{A finite modulus can still give an exact valuation}{18 / WHEN THE RESIDUE IS ENOUGH}{figures/18_modular_exactness.pdf}{
The valuation $v_2(S)$ counts the factors of $2$ in a nonzero integer $S$.

If $S\bmod2^K$ is nonzero, its trailing zero bits give the exact valuation below $K$.

A zero residue would only imply $v_2(S)\geq K$. The audit counts such unresolved residues explicitly.

For $K=128$ in the default run, none occurred. Every reported valuation in its finite domain is exact.
}{The word exact refers to the tested pairs. It does not extend the domain to infinity.}

\GuideSpread{A small independent scan reveals the texture}{19 / EXACT DIRECT SUMS}{figures/19_v2_heatmap.pdf}{
Every colored cell corresponds to one exact binomial sum with
\[
1\leq n\leq100,
\qquad0\leq m\leq4n.
\]
There are $20300$ such pairs. White cells are outside that domain.

This implementation uses direct sums and arbitrary precision integers. It checks the faster Pascal method over the same smaller range.

The uneven pattern is finite data. A uniform asymptotic valuation bound needs an argument beyond the image.
}{Independent arithmetic methods check that the visible pattern is produced correctly.}

\GuideSpread{The largest observed valuation grows in jumps}{20 / PREFIX MAXIMA}{figures/20_prefix_maxima.pdf}{
For each $N$, maximize $v_2(S_{n,m})$ over the triangular domain up to $N$.

The plot shows the independently computed prefixes through $N=100$. The maximum there is $20$.

The larger default scan, up to $N=2000$, reaches $31$.

Neither observation supplies a universal upper bound. Rare large valuations could occur at later indices, and an asymptotic claim requires control of them.
}{A record value summarizes a finite search. It does not control the unsearched tail.}

\GuideSpread{See the full domain of the large run}{21 / EIGHT MILLION PAIRS, WITH A BOUNDARY}{figures/21_scan_domain.pdf}{
The default scan tested
\[
1\leq n\leq2000,
\qquad0\leq m\leq4n.
\]
The total is
\[
2N^2+3N=8006000
\quad(N=2000).
\]
Its maximum valuation is $31$. The first maximizer in the scan's order is $(1466,2113)$.

The fresh run agrees exactly with the historical result. The unshaded region remains untested.
}{A large bounded search is strong implementation evidence and bounded mathematical evidence.}

\GuideSpread{The least common multiple carries a real cost}{22 / DENOMINATOR CLEARING}{figures/22_lcm_growth.pdf}{
Let
\[
L_t=\operatorname{lcm}(1,\ldots,t).
\]
The historical cleared pairs use $Q=8L_{2(n+m)}^2$.

The step curve uses exact integer $L_t$ through $t=300$. Its logarithm shows the scale of the clearing factor.

The line $t$ is a reference. The theorem $\log L_t\sim t$ is an external asymptotic input, equivalent to the prime number theorem in the standard formulation.
}{Clearing by an LCM converts rational coefficients to integers while multiplying the error.}

\GuideSpread{Five common divisors do not establish a rate}{23 / THE HISTORICAL CONTENT TABLE}{figures/23_gcd_rate.pdf}{
For a cleared integer pair, write
\[
d_{n,m}=\gcd(|A_{n,m}|,|B_{n,m}|).
\]
Dividing both coefficients by $d_{n,m}$ can improve the resulting integer form.

The five historical pairs were recomputed exactly. The chart displays $\log d_{n,m}/(n+m)$.

These five values decrease. That is a statement about this table, not a proof of an asymptotic rate along an infinite sequence.
}{The useful question is the sustained divisor growth in a specified normalization.}

\GuideSpread{Most of the large divisor comes from the chosen scale}{24 / MINIMAL CLEARING CHANGES THE INTERPRETATION}{figures/24_gcd_decomposition.pdf}{
Write the rational second packet coefficient as $H=N/D$ in lowest terms and the first as the integer $S$.

Minimal clearing gives $(DS,N)$. Historical clearing gives $R(DS,N)$, where $R=Q/D$ is an integer in these pairs.

Thus
\[
d_{\rm cleared}=R\,\gcd(DS,N).
\]
The five minimal pair gcds are $8,4,8,64,8$. A large cleared gcd can chiefly reflect a redundant clearing factor.
}{Separate recovered overclearing from intrinsic common divisibility.}

\GuideSpread{The missing inequality is a comparison of rates}{25 / A SCHEMATIC THRESHOLD MAP}{figures/25_threshold_plane.pdf}{
Suppose, along a chosen sequence,
\begin{align*}
|F_n|&=e^{-\delta n+o(n)},\\
Q_n&=e^{cn+o(n)},\\
d_n&=e^{hn+o(n)}.
\end{align*}
Then
\[
|Q_nF_n/d_n|
=e^{(c-\delta-h)n+o(n)}.
\]
Strictly positive margin $\delta+h-c$ gives decay. Equality needs finer estimates.

The figure fixes the illustrative value $\delta=1.5$. Its numerical boundary is not a Catalan theorem.
}{A successful proof needs a uniform estimate that crosses the boundary with a strict margin.}

\GuideSpread{A finite prefix cannot decide an infinite tail}{26 / COUNTERMODELS TO EXTRAPOLATION}{figures/26_finite_countermodels.pdf}{
The two drawn sequences agree through index $50$ and have deliberately different tails.

They are synthetic examples. They are not fitted predictions for Catalan coefficients.

Any bounded table can be extended in many incompatible ways unless a theorem restricts the extension.

An eventual rate therefore needs a uniform inequality, a structural identity or an external result with matching hypotheses.
}{Agreement on every observed point cannot replace a theorem about all later points.}

\GuideSpread{Even the sign deserves its own check}{27 / FILTERED NONVANISHING}{figures/27_filter_signs.pdf}{
The raw audit form is
\[
F_{n,m}=8(-1)^n(S_{n,m}G-H_{n,m}),
\]
where $H_{n,m}=\sum_i\binom mi16^{n+i}v_{n+i}$.

Exact interval arithmetic certifies nonzero signs in $860$ pairs through $n=20$.

The guess $\operatorname{sgn}F_{n,m}=(-1)^m$ fails at $(1,1)$:
\[
F_{1,1}=\frac{173056}{9}-20992G>0.
\]
Under the corrected source normalization, $F_{n,m}=16^n\widetilde J_{n,m}$. Even $m$ gives a positive integral. Odd $m$ has no comparable general sign rule here.
}{A failed simple sign rule is evidence that filtered nonvanishing needs a precise argument.}

\GuideSpread{An even filter keeps the integral positive}{27A / THE CORRECTED INTEGRAL}{figures/30_integral_filter.pdf}{
Put
\[
\begin{gathered}
\psi=\frac{xy(1-x)(1-y)}{1-xy},\\[2mm]
\dd\mu=\frac{\dd x\,\dd y}{\sqrt{x(1-y)}(1-xy)}.
\end{gathered}
\]
The corrected source identity is
\[
\widetilde J_{n,m}=\int_{[0,1]^2}
\psi^n(1-16\psi)^m\,\dd\mu.
\]
The square root does not cover $1-xy$. Also $J_0=8G$, which fixes the factor in $J_n=8(-1)^n(u_nG-v_n)$.

The range ends at $\phi^{-5}>1/16$, where $\phi=(1+\sqrt5)/2$. The filter changes sign inside the square.

For even $m$, the integrand is nonnegative and positive on an open set. Thus $\widetilde J_{n,m}>0$ for every $n\geq0$.
}{Choose even filter degrees to establish nonvanishing without extrapolating a sign table.}

\GuideSpread{Better rational approximation still falls short}{27B / THE SOURCE RATE BOUND}{figures/31_source_rate.pdf}{
Take even $m=m(n)$ with $m/n\to\alpha\geq0$. The source supplies
\[
\limsup_{n\to\infty}
|\widetilde J_{n,m}|^{1/n}\leq M(\alpha),
\]
where $M(\alpha)$ is the larger of the two branches displayed in the figure. The first has its continuous value $1/16$ at $\alpha=0$.

The reported reduced denominator exponent reaches about $0.62335$. It remains below the irrationality threshold $1$.

In the historical clearing, the exponential bound for the integer error is governed by
\[
\kappa(\alpha)=4(1+\alpha)+\log16+\log M(\alpha).
\]
This bound does not force decay. Positive $\kappa$ does not prove that the actual forms grow.

Further common divisibility or a sharper error estimate would need a uniform argument along infinitely many indices.
}{The source bound improves rational approximations. Its present arithmetic estimate does not establish irrationality.}

\GuideSpread{The conditional proof route is completely clear}{28 / WHAT WOULD FINISH THE ARGUMENT}{figures/28_conditional_route.pdf}{
Construct integer pairs $(A_n,B_n)$ such that the forms $A_nG-B_n$ are nonzero and tend to zero along an infinite sequence.

The elementary lattice argument then proves $G$ irrational.

The research problem is to meet those hypotheses simultaneously. The source bundle separates identities, denominator results, finite checks and unresolved estimates.

No finite scan establishes the necessary global decay of the normalized integer forms.
}{The endpoint is a conditional theorem. Applying it to Catalan's constant remains unfinished.}

\GuideSpread{Read the evidence with its proper scope}{29 / SOURCES AND REPRODUCTION}{figures/29_evidence_boundaries.pdf}{
The monograph gives proofs, notation translations and the detailed source contracts.

The evidence bundle preserves the historical audit and its outputs. It also contains the fresh run, independent diagnostics, figure source and figure provenance.

External sources include Zudilin's recurrence work and Eskandari's September 2026 preprint. The ledger records source versions, normalization repairs and open implications.

{\raggedright Reproduce the figures with \texttt{visual/figures.py}. Rounded logarithms and decimal coordinates are marked as display aids. Exact rational decisions remain in the computation files.\par}
}{The release records audited progress and its limits. No proof of Catalan's irrationality is established here.}
\end{document}
